%%%
%%% Radical "⽈"
%%%
\section*{Radical 73: ``⽈''}\addcontentsline{toc}{section}{Radical 73: ⽈}\addcontentsline{loh}{figure}{\#\#\#\# 73: ⽈}

%%%%%%%%%% 曰 %%%%%%%%%%
\subsection*{曰}\addcontentsline{loh}{figure}{曰}

\begin{Entry}{曰}{4}{⽈}
  \begin{Phonetics}{曰}{yue1}[][HSK 7-9]
    \definition{v.}{dizer | chamar; nomear | falar}
  \end{Phonetics}
\end{Entry}

%%%%%%%%%% 曲 %%%%%%%%%%
\subsection*{曲}\addcontentsline{loh}{figure}{曲}

\begin{Entry}{曲}{6}{⽈}
  \begin{Phonetics}{曲}{qu1}
    \definition*{s.}{Sobrenome: Qu}
    \definition{adj.}{dobrado; curvo; sinuoso | errado; injustificável}
    \definition{s.}{curva (de um rio, etc.) | fermento; levedura}
    \definition{v.}{dobrar; torcer}
  \antonymref{直}{zhi2}
  \end{Phonetics}
  \begin{Phonetics}{曲}{qu3}[][HSK 7-9]
    \definition{s.}{canção; melodia; partitura}
  \end{Phonetics}
\end{Entry}

\begin{Entry}{曲折}{6,7}{⽈,⼿}
  \begin{Phonetics}{曲折}{qu1zhe2}[][HSK 7-9]
    \definition{adj.}{sinuoso; tortuoso; não reto | complicado; intrincado; a situação e o enredo são complexos}
    \definition{s.}{complicações; enredo complexo e frustrante}
  \end{Phonetics}
\end{Entry}

\begin{Entry}{曲线}{6,8}{⽈,⽷}
  \begin{Phonetics}{曲线}{qu1xian4}[][HSK 7-9]
    \definition[条]{s.}{curva; em geometria, refere"-se à trajetória de um ponto que se move sob certas condições em um plano ou no espaço | corpo curvado; linhas onduladas; também se refere às linhas do corpo humano}
  \end{Phonetics}
\end{Entry}

\begin{Entry}{曲棍球}{6,12,11}{⽈,⽊,⽟}
  \begin{Phonetics}{曲棍球}{qu1gun4qiu2}
    \definition{s.}{hóquei em campo; hóquei | bola de hóquei}
  \end{Phonetics}
\end{Entry}

%%%%%%%%%% 曳 %%%%%%%%%%
\subsection*{曳}\addcontentsline{loh}{figure}{曳}

\begin{Entry}{曳}{6}{⽈}
  \begin{Phonetics}{曳}{ye4}
    \definition{v.}{arrastar; puxar; rebocar; sirgar; tracionar}
  \seealsoref{拽}{ye4}
  \end{Phonetics}
\end{Entry}

%%%%%%%%%% 更 %%%%%%%%%%
\subsection*{更}\addcontentsline{loh}{figure}{更}

\begin{Entry}{更}{7}{⽈}
  \begin{Phonetics}{更}{geng1}
    \definition*{s.}{Sobrenome: Geng}
    \definition{clas.}{um dos cinco períodos de duas horas em que a noite era anteriormente dividida; vigília; antigamente, a noite era dividida em cinco turnos, cada um com aproximadamente duas horas de duração}
    \definition{v.}{alterar; substituir | experimentar}
  \seealsoref{更加}{geng4jia1}
  \synonymref{改}{gai3}
  \synonymref{更加}{geng4jia1}
  \end{Phonetics}
  \begin{Phonetics}{更}{geng4}[][HSK 2]
    \definition{adv.}{mais; ainda mais | além disso; além do mais; ainda mais}
  \end{Phonetics}
\end{Entry}

\begin{Entry}{更加}{7,5}{⽈,⼒}
  \begin{Phonetics}{更加}{geng4jia1}[][HSK 3]
    \definition{adv.}{mais; ainda mais; em maior grau; indica um nível mais profundo ou um aumento ou diminuição quantitativa adicional}
  \end{Phonetics}
\end{Entry}

\begin{Entry}{更衣室}{7,6,9}{⽈,⾐,⼧}
  \begin{Phonetics}{更衣室}{geng1yi1shi4}[][HSK 7-9]
    \definition{s.}{vestiário | camarim | \emph{toilet}; toucador}
  \end{Phonetics}
\end{Entry}

\begin{Entry}{更改}{7,7}{⽈,⽁}
  \begin{Phonetics}{更改}{geng1gai3}[][HSK 7-9]
    \definition{v.}{alterar; mudar}
  \end{Phonetics}
\end{Entry}

\begin{Entry}{更是}{7,9}{⽈,⽇}
  \begin{Phonetics}{更是}{geng4shi4}[][HSK 6]
    \definition{adv.}{ainda mais (assim)}
  \end{Phonetics}
\end{Entry}

\begin{Entry}{更换}{7,10}{⽈,⼿}
  \begin{Phonetics}{更换}{geng1huan4}[][HSK 5]
    \definition{v.}{alterar; mudar; substituir; comutar}
  \end{Phonetics}
\end{Entry}

\begin{Entry}{更新}{7,13}{⽈,⽄}
  \begin{Phonetics}{更新}{geng1xin1}[][HSK 5]
    \definition{v.}{renovar; atualizar; substituir; remover o antigo e substituir pelo novo}
  \end{Phonetics}
\end{Entry}

%%%%%%%%%% 曹 %%%%%%%%%%
\subsection*{曹}\addcontentsline{loh}{figure}{曹}

\begin{Entry}{曹}{11}{⽈}
  \begin{Phonetics}{曹}{cao2}
    \definition*{s.}{Cao, um estado vassalo na Dinastia Zhou | Sobrenome: Cao}
    \definition{s.}{Literário: pessoas do mesmo tipo | Arcaico: departamento do governo sob o monarca}
  \end{Phonetics}
\end{Entry}

\begin{Entry}{曹雪芹}{11,11,7}{⽈,⾬,⾋}
  \begin{Phonetics}{曹雪芹}{cao2 xue3qin2}
    \definition*{s.}{Cao Xueqin (c. 1715--c. 1764), autor reconhecido de O Sonho da Câmara Vermelha 《红楼梦》}
  \seealsoref{红楼梦}{hong2lou2 meng4}
  \end{Phonetics}
\end{Entry}

%%%%%%%%%% 曾 %%%%%%%%%%
\subsection*{曾}\addcontentsline{loh}{figure}{曾}

\begin{Entry}{曾}{12}{⽈}
  \begin{Phonetics}{曾}{ceng2}[][HSK 4]
    \definition{adv.}{indica que uma ação já aconteceu ou um estado já existiu}
  \end{Phonetics}
  \begin{Phonetics}{曾}{zeng1}
    \definition*{s.}{Sobrenome: Zeng}
    \definition{s.}{relacionamento entre bisnetos e bisavós; (parentesco) duas gerações de diferença}
  \end{Phonetics}
\end{Entry}

\begin{Entry}{曾经}{12,8}{⽈,⽷}
  \begin{Phonetics}{曾经}{ceng2jing1}[][HSK 3]
    \definition{adv.}{uma vez; indica que houve algum comportamento ou situação}
  \end{Phonetics}
\end{Entry}

%%%%%%%%%% 替 %%%%%%%%%%
\subsection*{替}\addcontentsline{loh}{figure}{替}

\begin{Entry}{替}{12}{⽈}
  \begin{Phonetics}{替}{ti4}[][HSK 4]
    \definition{prep.}{para; em nome de}
    \definition{s.}{decadência; declínio; enfraquecimento}
    \definition{v.}{substituir; substituir por; tomar o lugar de}
  \end{Phonetics}
\end{Entry}

\begin{Entry}{替代}{12,5}{⽈,⼈}
  \begin{Phonetics}{替代}{ti4dai4}[][HSK 4]
    \definition{v.}{substituir; suplantar}
  \end{Phonetics}
\end{Entry}

\begin{Entry}{替身}{12,7}{⽈,⾝}
  \begin{Phonetics}{替身}{ti4shen1}[][HSK 7-9]
    \definition{s.}{bode expiatório | dublê de corpo | dublê}
    \definition{v.}{substituir; substituir alguém; ocupar o lugar de}
  \end{Phonetics}
\end{Entry}

\begin{Entry}{替换}{12,10}{⽈,⼿}
  \begin{Phonetics}{替换}{ti4huan4}[][HSK 7-9]
    \definition{v.}{substituir; substituir por; deslocar; tomar o lugar de; alternar}
  \synonymref{撤换}{che4huan4}
  \synonymref{更换}{geng1huan4}
  \synonymref{交换}{jiao1huan4}
  \synonymref{替代}{ti4dai4}
  \antonymref{代替}{dai4ti4}
  \end{Phonetics}
\end{Entry}

%%%%%%%%%% 最 %%%%%%%%%%
\subsection*{最}\addcontentsline{loh}{figure}{最}

\begin{Entry}{最}{12}{⽈}
  \begin{Phonetics}{最}{zui4}[][HSK 1]
    \definition{adv.}{(diante de um adjetivo ou verbo) o mais | (colocado antes de um substantivo de localidade ou de uma palavra que indica um lugar)  mais distante ou mais próximo de (um lugar) | mais; melhor; pior; primeiro; muito; menos; acima de tudo; indica que uma determinada característica excede todas as outras pessoas ou coisas do mesmo tipo}
    \definition{s.}{o máximo; o melhor (ou o mais alto, o maior, etc.)}
  \synonymref{极}{ji2}
  \end{Phonetics}
\end{Entry}

\begin{Entry}{最少}{12,4}{⽈,⼩}
  \begin{Phonetics}{最少}{zui4shao3}
    \definition{adj.}{menor; menos; mínimo; indica o limite mínimo}
    \definition{adv.}{pelo menos; no mínimo}
  \seealsoref{至少}{zhi4shao3}
  \synonymref{起码}{qi3ma3}
  \synonymref{至少}{zhi4shao3}
  \antonymref{最多}{zui4 duo1}
  \end{Phonetics}
\end{Entry}

\begin{Entry}{最优}{12,6}{⽈,⼈}
  \begin{Phonetics}{最优}{zui4you1}
    \definition{expr.}{ótimo}
  \end{Phonetics}
\end{Entry}

\begin{Entry}{最先}{12,6}{⽈,⼉}
  \begin{Phonetics}{最先}{zui4xian1}
    \definition{adv.}{(o) primeiro}
  \end{Phonetics}
\end{Entry}

\begin{Entry}{最后}{12,6}{⽈,⼝}
  \begin{Phonetics}{最后}{zui4hou4}[][HSK 1]
    \definition{s.}{último; final; definitivo; refere"-se ao tempo, local, etc. que vem depois de outros tempos, locais, etc. na ordem sequencial}
  \seealsoref{结果}{jie2/guo3}
  \seealsoref{终于}{zhong1yu2}
  \synonymref{结果}{jie2/guo3}
  \synonymref{结尾}{jie2wei3}
  \synonymref{期末}{qi1mo4}
  \synonymref{终于}{zhong1yu2}
  \synonymref{最终}{zui4zhong1}
  \antonymref{初步}{chu1bu4}
  \antonymref{起初}{qi3chu1}
  \antonymref{首先}{shou3xian1}
  \antonymref{率先}{shuai4xian1}
  \antonymref{最初}{zui4chu1}
  \end{Phonetics}
\end{Entry}

\begin{Entry}{最多}{12,6}{⽈,⼣}
  \begin{Phonetics}{最多}{zui4 duo1}
    \definition{adj.}{mais; máximo}
    \definition{adv.}{no máximo}
  \antonymref{最少}{zui4shao3}
  \end{Phonetics}
\end{Entry}

\begin{Entry}{最好}{12,6}{⽈,⼥}
  \begin{Phonetics}{最好}{zui4hao3}[][HSK 1]
    \definition{adj.}{melhor; de primeira qualidade; excelente}
    \definition{adv.}{seria melhor; seria o ideal; indica a escolha mais adequada entre várias possibilidades}
  \synonymref{顶级}{ding3ji2}
  \synonymref{最佳}{zui4jia1}
  \end{Phonetics}
\end{Entry}

\begin{Entry}{最初}{12,7}{⽈,⾐}
  \begin{Phonetics}{最初}{zui4chu1}[][HSK 4]
    \definition{adj.}{primordial; inicial; primeiro}
    \definition{adv.}{inicialmente; originalmente}
    \definition{s.}{o período mais antigo; início; começo}
  \end{Phonetics}
\end{Entry}

\begin{Entry}{最近}{12,7}{⽈,⾡}
  \begin{Phonetics}{最近}{zui4jin4}[][HSK 2]
    \definition{adj.}{mais próximo}
    \definition{s.}{recentemente; ultimamente; de tarde; refere"-se aos dias antes ou logo depois de um discurso | em breve; no futuro próximo; o futuro próximo}
  \end{Phonetics}
\end{Entry}

\begin{Entry}{最远}{12,7}{⽈,⾡}
  \begin{Phonetics}{最远}{zui4yuan3}
    \definition{adv.}{à distância máxima | mais distante}
  \end{Phonetics}
\end{Entry}

\begin{Entry}{最佳}{12,8}{⽈,⼈}
  \begin{Phonetics}{最佳}{zui4jia1}[][HSK 6]
    \definition{adj.}{o melhor (atleta, filme etc); ótimo}
  \end{Phonetics}
\end{Entry}

\begin{Entry}{最终}{12,8}{⽈,⽷}
  \begin{Phonetics}{最终}{zui4zhong1}[][HSK 6]
    \definition{adv.}{finalmente; por fim}
    \definition{s.}{final; definitivo}
  \end{Phonetics}
\end{Entry}

\begin{Entry}{最高}{12,10}{⽈,⾼}
  \begin{Phonetics}{最高}{zui4gao1}
    \definition{adj.}{o mais alto; supremo; mais alto}
    \definition[类]{s.}{teto de acumulação; taxa máxima de acumulação}
  \synonymref{极限}{ji2xian4}
  \end{Phonetics}
\end{Entry}

\begin{Entry}{最善}{12,12}{⽈,⼝}
  \begin{Phonetics}{最善}{zui4shan4}
    \definition{adj.}{ideal | o melhor}
  \end{Phonetics}
\end{Entry}

\begin{Entry}{最新}{12,13}{⽈,⽄}
  \begin{Phonetics}{最新}{zui4xin1}
    \definition{adv.}{mais recente; atualizado; em dia | mais novo}
  \synonymref{前沿}{qian2yan2}
  \synonymref{新潮}{xin1chao2}
  \synonymref{新式}{xin1shi4}
  \synonymref{新鲜}{xin1xian1}
  \synonymref{崭新}{zhan3xin1}
  \antonymref{陈旧}{chen2jiu4}
  \antonymref{古老}{gu3lao3}
  \antonymref{过时}{guo4shi2}
  \antonymref{落后}{luo4/hou4}
  \end{Phonetics}
\end{Entry}

%%%%% EOF %%%%%

